In this section, we describe the details of the IFT setting of \citet{xia2024less},
as well as the details of our method.

\paragraph{Setting.}
The setting contains a fixed data pool: instruction fine-tuning data from a data
pool consisting of four combined IFT datasets (cf. Table~\ref{tab:ift_training_dataset}
and \citet{xia2024less} for more information). The goal is to select the data
that yields the best possible task performance for a LoRA fine-tuning run.
We adapt a LoRA to a Gemma-2B model (the pretraining-only Gemma-2B model)
using the LoRA configuration from \citet{xia2024less}.

\paragraph{Data splits.}
See Table~\ref{tab:datasets} for a description of the available data for each
task, along with the task setup details. \citet{xia2024less} constructed these
extra samples by drawing from the ICL samples given in the tasks originally.
Note that we drop TydiQA from the original work of \citet{xia2024less} as there
are not enough samples to select with (there is only one from each category, for
a total of 7). 

\paragraph{Method.}
We execute Algorithm~\ref{alg:depsing} with $k$ as 150 steps from the end of training and
the Bernoulli parameter $p$ controlling the step size as 0.2. At each step, we
choose a ``minibatch'' with a size equal to half the target set and a quarter of
the target set for BBH and MMLU, respectively (that is, we only select to
optimize performance on a fraction of the target set at a time). We model
select over iterates and hyperparameters by (a) choosing the top three steps in
terms of validation loss for each run (b) selecting the best one in terms of
full train set accuracy (including the part that we trained on). We perform this
procedure---akin to Pareto optimization~\citep{jin2008pareto}---because the
validation set is so small (as the overall set of samples is very small) that it
is difficult to select models without overfitting otherwise.

We compare with two baselines: training on the full dataset (i.e., training on
the entirety of all the data for a single epoch), and LESS (we use the data
selected according to ``LESS-T''~\citep{xia2024less}, following the
recommendation of 4 epochs).

For model training, we train with ADAM ($\beta_1=0.95, \beta_2=0.975$,
decoupled weight decay as $10^{-5}$) and a one-cycle linear schedule starting
at $10^{-6}$ of the maximum learning rate, reaching the peak over 25\% of
training, then ending at 0.1 of the maximum learning rate. We insert a positive
$\epsilon_\textrm{root}$ into the inverse square root term in the ADAM update
to prevent metagradient (and to a lesser extent update) blowup (see Eq.\
\ref{eq:adam_update}). The model training is the same across selected data,
except that we use $\epsilon_\textrm{root}=10^{-7}$ for MGD-selected data and
$\epsilon_\textrm{root}=10^{-9}$ for the other runs (we select the optimal
parameter for each class of method). We additionally hyperparameter select for
the best learning rate across each baseline by minimizing validation set loss;
LESS performs best with a smaller learning rate ($0.00024$ for BBH and
$0.00012$ for MMLU) than training on the full dataset or with MGD ($0.0006$ for
both). We normalize the loss of each training sample by taking the mean across
predicted tokens during training, and do not divide by the batch size
(important for scaling the $\epsilon_\mathrm{root}$ term, but otherwise ADAM is
invariant to the scale). 

\paragraph{Selecting smooth model training for MGD.} For MGD runs, we jointly
select learning rate and $\epsilon_{\mathrm{root}}$ using the smoothness metric
of Section~\ref{sec:localpred}. We find that the choice of $\epsilon_\mathrm{root}$ term
is important (just as the choice of $\epsilon$ is important in standard ADAM training); choosing a much larger term results in non-smooth training. We also find that
metagradients are sensitive to learning rate schedule; choosing a much larger or
smaller maximum learning rate results in non-smooth training.

\begin{table}[h]
    \caption{Overview of datasets used in IFT dataset selection (from \citet{xia2024less}).}
    \label{tab:datasets}
    \centering
    \begin{tabular}{lccccccc}
        \toprule
        \textbf{Dataset} & \textbf{\# Shot} & \textbf{\# Tasks} & $n_\mathrm{target}$ & $n_\mathrm{val}$  & $n_\mathrm{test}$ & \textbf{Answer Type} & \textbf{Type of Task} \\
        \midrule
        MMLU & 5 & 57 & 57 & 228 & 18,721 & Letter options & Knowledge/Recall \\
        BBH & 3 & 23 & 23 & 46 & 920 & COT and answer & Reasoning \\
        \bottomrule
    \end{tabular}
\end{table}

\begin{table}[h]
    \centering
    \caption{Details of IFT training datasets.}
    \label{tab:ift_training_dataset}
    \begin{tabular}{l l l c c}
        \toprule
        Dataset & \# Instance & Sourced from & Prompt Len. & Completion Len. \\
        \midrule
        FLAN V2 & 100,000 & NLP datasets and human-written instructions & 355.7 & 31.2 \\
        CoT & 100,000 & NLP datasets and human-written CoTs & 266 & 53.2 \\
        Dolly & 15,011 & Human-written from scratch & 118.1 & 91.3 \\
        Open Assistant 1 & 55,668 & Human-written from scratch & 34.8 & 212.5 \\
        \bottomrule
    \end{tabular}
\end{table}

\paragraph{IFT results} \quad %

\begin{figure}[hb]
	\centering
	\caption{MGD dataset selection improves the validation loss over metagradient steps, demonstrating our method's efficacy.  However, the gap between loss on samples MGD directly optimizes on and the validation samples widens over the number of iterates, and there is overfitting depending on the number of steps taken. }
	\includegraphics[width=0.45\textwidth]{plots/bbh_loss_over_time.pdf}
	\quad
	\includegraphics[width=0.45\textwidth]{plots/mmlu_loss_over_time.pdf}
	\label{fig:bbh_mmlu_over_time}
\end{figure}
